% Diagramas de comparacion de pipelines
% Para usar en el capitulo 5 de la tesis (Aislamiento del Modulo TLR)

% IMPORTANTE: Desactivar shorthand de babel para que TikZ funcione
\shorthandoff{<>}

% ============================================================
% DIAGRAMA 1: PIPELINE APOLLO ORIGINAL
% ============================================================

\begin{figure}[H]
\centering
\begin{tikzpicture}[
    node distance=0.8cm,
    box/.style={rectangle, draw, rounded corners, minimum width=3cm, minimum height=0.8cm, align=center, font=\small},
    inputbox/.style={box, fill=blue!15},
    procbox/.style={box, fill=orange!20},
    outputbox/.style={box, fill=green!15},
    arrow/.style={->, thick, >=stealth}
]

% Inputs
\node[inputbox] (img) {Imagen de Camara};
\node[inputbox, right=0.5cm of img] (hdmap) {HD-Map\\(Coords 3D)};
\node[inputbox, right=0.5cm of hdmap] (pose) {Pose Vehiculo\\(GPS + TF)};
\node[inputbox, right=0.5cm of pose] (calib) {Calibracion\\Camara};

% Etapas
\node[procbox, below=1.8cm of hdmap, xshift=0.5cm] (prep) {\textbf{Region Proposal}\\Query HD-Map, Proyeccion 3D a 2D};

\node[procbox, below=0.8cm of prep] (detect) {\textbf{Detection}\\CNN SSD-style + NMS};

\node[procbox, below=0.8cm of detect] (select) {\textbf{Assignment}\\Algoritmo Hungaro\\(70\% dist, 30\% conf)};

\node[procbox, below=0.8cm of select] (recog) {\textbf{Recognition}\\3 Clasificadores CNN\\(vert/quad/hori)};

\node[procbox, below=0.8cm of recog] (track) {\textbf{Tracking}\\Reglas temporales\\+ Semantic Voting};

% Output
\node[outputbox, below=0.8cm of track] (output) {\textbf{Output}\\Color + Blink + Signal ID};

% Coordenada para la linea horizontal comun
\coordinate (linea_comun) at ($(img.south)+(0,-0.5)$);

% Arrows inputs - todas bajan a la misma altura y luego van al centro
\draw[arrow] (img.south) -- (img.south |- linea_comun) -| (prep.north);
\draw[arrow] (hdmap.south) -- (hdmap.south |- linea_comun) -- (prep.north |- linea_comun) -- (prep.north);
\draw[arrow] (pose.south) -- (pose.south |- linea_comun) -| (prep.north);
\draw[arrow] (calib.south) -- (calib.south |- linea_comun) -| (prep.north);

% Arrows between stages
\draw[arrow] (prep) -- (detect);
\draw[arrow] (detect) -- (select);
\draw[arrow] (select) -- (recog);
\draw[arrow] (recog) -- (track);
\draw[arrow] (track) -- (output);

% Modulos Apollo
\node[right=0.3cm of prep, font=\tiny\ttfamily, text=gray] {traffic\_light\_region\_proposal};
\node[right=0.3cm of detect, font=\tiny\ttfamily, text=gray] {traffic\_light\_detection};
\node[right=0.3cm of select, font=\tiny\ttfamily, text=gray] {traffic\_light\_detection};
\node[right=0.3cm of recog, font=\tiny\ttfamily, text=gray] {traffic\_light\_recognition};
\node[right=0.3cm of track, font=\tiny\ttfamily, text=gray] {traffic\_light\_tracking};

\end{tikzpicture}
\caption{Pipeline de Apollo original con los 4 modulos}
\label{fig:pipeline_apollo}
\end{figure}


% ============================================================
% DIAGRAMA 2: PIPELINE REIMPLEMENTACION
% ============================================================

\begin{figure}[H]
\centering
\begin{tikzpicture}[
    node distance=0.8cm,
    box/.style={rectangle, draw, rounded corners, minimum width=3cm, minimum height=0.8cm, align=center, font=\small},
    inputbox/.style={box, fill=blue!15},
    procbox/.style={box, fill=orange!20},
    outputbox/.style={box, fill=green!15},
    simplebox/.style={box, fill=yellow!20, dashed},
    arrow/.style={->, thick, >=stealth}
]

% Inputs (simplificados)
\node[inputbox] (img) {Imagen de Camara};
\node[simplebox, right=1cm of img] (boxes) {Projection Boxes\\(archivo .txt)};

% Etapas
\node[procbox, below=1cm of img, xshift=1.5cm] (prep) {\textbf{Preprocessing}\\Crop + Resize 270x270};

\node[procbox, below=0.8cm of prep] (detect) {\textbf{Detection}\\CNN SSD-style + NMS};

\node[procbox, below=0.8cm of detect] (select) {\textbf{Assignment}\\Algoritmo Hungaro\\(70\% dist, 30\% conf)};

\node[procbox, below=0.8cm of select] (recog) {\textbf{Recognition}\\3 Clasificadores CNN\\(vert/quad/hori)};

\node[procbox, below=0.8cm of recog] (track) {\textbf{Tracking}\\Reglas temporales};

% Output
\node[outputbox, below=0.8cm of track] (output) {\textbf{Output}\\Color + Blink + Signal ID};

% Arrows inputs
\draw[arrow] (img.south) -- ++(0,-0.3) -| (prep.north);
\draw[arrow] (boxes.south) -- ++(0,-0.3) -| (prep.north);

% Arrows between stages
\draw[arrow] (prep) -- (detect);
\draw[arrow] (detect) -- (select);
\draw[arrow] (select) -- (recog);
\draw[arrow] (recog) -- (track);
\draw[arrow] (track) -- (output);

% Archivos Python
\node[right=0.3cm of prep, font=\tiny\ttfamily, text=gray] {utils.py};
\node[right=0.3cm of detect, font=\tiny\ttfamily, text=gray] {detector.py};
\node[right=0.3cm of select, font=\tiny\ttfamily, text=gray] {selector.py};
\node[right=0.3cm of recog, font=\tiny\ttfamily, text=gray] {recognizer.py};
\node[right=0.3cm of track, font=\tiny\ttfamily, text=gray] {tracking.py};

% Nota de simplificacion
\node[below=0.3cm of output, font=\footnotesize, text=gray, align=center] {Linea punteada = simplificacion\\respecto a Apollo};

\end{tikzpicture}
\caption{Pipeline del modulo aislado (con interfaz simplificada)}
\label{fig:pipeline_aislado}
\end{figure}


% ============================================================
% DIAGRAMA 3: COMPARACION LADO A LADO
% ============================================================

\begin{figure}[H]
\centering
\begin{tikzpicture}[
    node distance=0.6cm,
    box/.style={rectangle, draw, rounded corners, minimum width=2.2cm, minimum height=0.6cm, align=center, font=\scriptsize},
    inputbox/.style={box, fill=blue!15},
    procbox/.style={box, fill=orange!20},
    outputbox/.style={box, fill=green!15},
    simplebox/.style={box, fill=yellow!20, dashed},
    greybox/.style={box, fill=gray!30},
    arrow/.style={->, thick, >=stealth},
    eqarrow/.style={<->, thick, >=stealth, color=blue!60}
]

% ========== APOLLO (izquierda) ==========
\node[font=\bfseries] (titleA) {Apollo Original};

\node[inputbox, below=0.3cm of titleA, xshift=-2cm] (imgA) {Imagen};
\node[inputbox, right=0.5cm of imgA] (hdmapA) {HD-Map};
\node[inputbox, right=0.5cm of hdmapA] (poseA) {Pose};

\node[procbox, below=1.2cm of hdmapA] (prepA) {Region Proposal\\(3D a 2D)};
\node[procbox, below=0.5cm of prepA] (detectA) {Detection};
\node[procbox, below=0.5cm of detectA] (selectA) {Assignment};
\node[procbox, below=0.5cm of selectA] (recogA) {Recognition};
\node[procbox, below=0.5cm of recogA] (trackA) {Tracking};
\node[outputbox, below=0.5cm of trackA] (outA) {Output};

% Coordenada para linea horizontal Apollo
\coordinate (linea_A) at ($(imgA.south)+(0,-0.3)$);

% Arrows Apollo - todas a la misma altura
\draw[arrow] (imgA.south) -- (imgA.south |- linea_A) -| (prepA.north);
\draw[arrow] (hdmapA.south) -- (hdmapA.south |- linea_A) -- (prepA.north |- linea_A) -- (prepA.north);
\draw[arrow] (poseA.south) -- (poseA.south |- linea_A) -| (prepA.north);
\draw[arrow] (prepA) -- (detectA);
\draw[arrow] (detectA) -- (selectA);
\draw[arrow] (selectA) -- (recogA);
\draw[arrow] (recogA) -- (trackA);
\draw[arrow] (trackA) -- (outA);

% ========== REIMPLEMENTACION (derecha) ==========
\node[font=\bfseries, right=5cm of titleA] (titleR) {Modulo Aislado};

\node[inputbox, below=0.3cm of titleR, xshift=-0.8cm] (imgR) {Imagen};
\node[simplebox, right=0.5cm of imgR] (boxesR) {Boxes (txt)};

\node[simplebox, below=1.2cm of imgR, xshift=0.8cm] (prepR) {Preprocessing\\(2D directo)};
\node[procbox, below=0.5cm of prepR] (detectR) {Detection};
\node[procbox, below=0.5cm of detectR] (selectR) {Assignment};
\node[procbox, below=0.5cm of selectR] (recogR) {Recognition};
\node[procbox, below=0.5cm of recogR] (trackR) {Tracking};
\node[outputbox, below=0.5cm of trackR] (outR) {Output};

% Coordenada para linea horizontal Reimpl
\coordinate (linea_R) at ($(imgR.south)+(0,-0.3)$);

% Arrows Reimpl - todas a la misma altura
\draw[arrow] (imgR.south) -- (imgR.south |- linea_R) -| (prepR.north);
\draw[arrow] (boxesR.south) -- (boxesR.south |- linea_R) -| (prepR.north);
\draw[arrow] (prepR) -- (detectR);
\draw[arrow] (detectR) -- (selectR);
\draw[arrow] (selectR) -- (recogR);
\draw[arrow] (recogR) -- (trackR);
\draw[arrow] (trackR) -- (outR);

% ========== EQUIVALENCIAS ==========
\draw[eqarrow] (detectA.east) -- node[above, font=\tiny] {=} (detectR.west);
\draw[eqarrow] (selectA.east) -- node[above, font=\tiny] {=} (selectR.west);
\draw[eqarrow] (recogA.east) -- node[above, font=\tiny] {=} (recogR.west);
\draw[eqarrow] (trackA.east) -- node[above, font=\tiny] {=} (trackR.west);

% Diferencia
\draw[dashed, red, thick] (prepA.east) -- node[above, font=\tiny, red] {$\neq$} (prepR.west);

% Leyenda
\node[below=0.8cm of outA, xshift=2cm, font=\scriptsize, align=left] {
    \textcolor{blue}{$\leftrightarrow$ Equivalente funcionalmente}\\
    \textcolor{red}{- - - Simplificado (no HD-Map)}
};

\end{tikzpicture}
\caption[Comparación lado a lado: Apollo vs módulo aislado.]{Comparacion lado a lado: Apollo original vs modulo aislado. Las flechas azules indican etapas funcionalmente equivalentes; la linea roja punteada indica la unica etapa cuya interfaz se simplifico (HD-Map dinamico reemplazado por projection boxes precalculadas).}
\label{fig:pipeline_comparison}
\end{figure}

% Reactivar shorthand de babel
\shorthandon{<>}